\scanpage{217}
\deleted{Donc si $\beta(\text{classe }D)=0$, \unsure{alors|donc}}

Donc le noyau de $\beta$ [qui s'identifie à $H^1(X,\underline{\mathcal O}_X^*)\unsure{-|\to}H^1(U,\underline{\mathcal O}_U^*)$] est le sous-groupe de $H^1(X,\mathcal O_X^*)$ engendré par les classes de $D_1$ et $D_2$, donc par celle de $D_1$, puis que $D_1+D_2\sim0$. \deleted{Donc} Or $\operatorname{cl}(D_1)$ n'est pas un élément de torsion de $H^1(X,\underline{\mathcal O}_X^*)$, grâce à lemme d). Se servant du corollaire 1 \inserted{et du lemme c}, on trouve aussitôt le cor. 2.

\uline{Lemme e} Sous les conditions du lemme \unsure{b|s}, Corollaire 3, l'hom $(*)$ est \uline{bijectif}.

Il est en effet surjectif grâce au corollaire précédent.

On a ainsi \unsure{complètement|complété} \unsure{le|de|la} foncteur de Picard associé à $f:X\to S$.

\uline{Lemme f} Le foncteur \inserted{$\unsure{F|\mathscr F|\mathcal F}$} de Picard \deleted{\unsure{n'}}est \deleted{pas} représentable par un préschéma sur $S$, loc de type fini.

\deleted{Considérons en effet la fibre $f^{-1}(a)$, et l'élément de degrés $(1,\unsure{-1|1})$ [correspondant à $L_2$] de $H^1(f^{-1}(a),\underline{\mathcal O}_{f^{-1}(a)}^*)=F(\operatorname{Spec} k(a))$.}
% archive sheet number: 216
